% Carrega o pacote de cores, que deve vir ANTES do listings
\usepackage{xcolor}

\usepackage{microtype}

% Carrega o pacote de listings (ESSA LINHA ESTAVA FALTANDO)
\usepackage{listings}

% Agora sim, você pode configurar o pacote
\lstset{
    basicstyle=\ttfamily,
    numbers=left,
    keywordstyle=\color[rgb]{0.13,0.29,0.53}\bfseries,
    stringstyle=\color[rgb]{0.31,0.60,0.02},
    commentstyle=\color[rgb]{0.56,0.35,0.01}\itshape,
    numberstyle=\footnotesize,
    stepnumber=1,
    numbersep=5pt,
    backgroundcolor=\color[RGB]{248,248,248},
    showspaces=false,
    showstringspaces=false,
    showtabs=false,
    tabsize=2,
    captionpos=b,
    breaklines=true,         % <-- MANTENHA ESTA LINHA
    breakautoindent=true,
    escapeinside={\%*}{*)},
    linewidth=\textwidth,    % <-- Veja a dica bônus abaixo
    basewidth=0.5em,
}
